\section{Obstructions, limits and pitfalls}\label{sec:obstructions}

\paragraph{What this section is for (not used in the proof).}
This section records what did not work, where the method stops, and the errors
the proof had to avoid. Nothing in it is a premise of the main theorem, and no
proof of Parts~II--V cites it. Each statement is quoted, with its scope, from
the place where it is derived (a section of this paper, the archived manuscripts
of \cref{subsec:cert:archive}, or a research note), or is checked here by an
elementary argument (\cref{lem:obs:rejection-covers}, \cref{rem:obs:root-support}(i))
or a stated enumeration.
It serves a reader
who wants to know why the construction of \cref{sec:design} has its shape, and
one who wants to push the exponent further. \Cref{sec:history} tells the
development in order; here the material is grouped by kind: the accounting
(\cref{subsec:obs:plain,subsec:obs:reserved,subsec:obs:pairing}), the prime
ranges (\cref{subsec:obs:prefix,subsec:obs:class-one}), the root law and twelve
queries (\cref{subsec:obs:root,subsec:obs:twelve}), the obstruction
polynomials of the earlier recurrence (\cref{subsec:obs:polynomials}), limits
(\cref{subsec:obs:limits}), tested directions (\cref{subsec:obs:tested}),
pitfalls (\cref{subsec:obs:pitfalls}) and the scope of the reviews
(\cref{subsec:obs:reviews}). As in \cref{sec:history}, ``row~$n$'' is a row of
\cref{tab:history:checkpoints}, ``old \S$n$'' a section of the archived
manuscript of $644$ pages and ``master \S$n$'' a section of the archived
manuscript of $770$ pages, whose labels are quoted as master labels. Numbers
marked ``about'' are not certified.

\subsection{The crude bound exceeds \texorpdfstring{$1$}{1} (distinct from \texorpdfstring{$\Theta<1$}{Theta<1})}\label{subsec:obs:plain}

Without classes, without the common mean and with the norm envelopes as
matching envelopes (so that the dominator price is
$T=S+A$, the symmetric plus skew price, by the norm-envelope
case of (CP6) in \cref{def:cp6}), the comparison argument of
\cref{sec:comparison-proof} gives only the \emph{plain criterion}
$L_0+\mathcal E+J+W<1$, with $\mathcal E=B(e^S\cosh A-1)$ and
$J=Be^{S+A}(R_I/2)^6$. Theorem~C (\cref{thm:comparison}) refines it in three
independent ways, described in \cref{subsec:walk:accounting}: the classes of
(CP5) (\cref{sec:variance,sec:envelopes}), the direct Gram
(\cref{sec:matching-mean}), and the common mean (\cref{lem:record:common-mean}),
with which the trivial credit $m=1-L_0$ gives the closing
$L_0+\frac{63}{64}\mathcal E+J+W$ and a credit $m>1$ lowers it by a further
$(m-1+L_0)/64$ (\cref{rem:cp:trivial-mean}).

\begin{remark}[Why the plain criterion is refined]\label{rem:record:plain}
At the October~18 checkpoint ($\eps=441/44375000$,
\cref{subsec:history:layers}) the plain criterion fails even when the reserved
cost already contains the variance discounts of degrees one and two: its value
is then $1.1219\ldots>1$. The certified ledger of that checkpoint, which its
record verifier checks (\cref{subsec:cert:archive}), shows what each refinement
bought (upper bounds, rounded up; the gains are differences of the displayed
values):
\begin{center}
\small
\begin{tabular}{@{}>{\raggedright\arraybackslash}p{0.62\textwidth}rr@{}}
\toprule
After this step & Closing & Gain\\
\midrule
Variance discounts of degrees one and two; plain matching term & $1.121918$ & \\
Direct old/new $F$ matching Gram & $1.012226$ & $0.109692$\\
Variance discounts of degrees three and four & $1.012152$ & $0.000074$\\
Common mean, reserved coefficient $63/64$ & $0.999627$ & $0.012525$\\
Positive $F$ singleton mean, doubled & $0.995618$ & $0.004009$\\
Distinct-block pair mean and its polarized copy & $0.995490$ & $0.000128$\\
Extra direct-envelope singleton and pair means & $0.995315$ & $0.000175$\\
\bottomrule
\end{tabular}
\end{center}
Each line is a separate valid inequality between the same quantities, and no saving is counted in two lines.
\end{remark}

For the baseline construction \cref{subsec:walk:accounting} gives the same
decomposition at the values of \cref{tab:closing-values}: there the plain
criterion is about $1.603$, the closing without the common mean and the mean
credit about $1.021$, and every refinement is needed. The accounting also has
limits. The
coefficient $63/64$ was also derived pattern by pattern; the stronger
coefficient $31/32$ of that derivation needs the comparison and the matching to
use one bracket for every index-support pattern, which fails once the direct Gram prices the
$F$ records. At the October~18 checkpoint the positive mean of the
reserved-diagonal records had the ceiling $(e^{U_0}-1)/64$, about $0.016$, and
the coupons took about $0.0043$ of it. The final mean credit saves about
$0.0092$: its coupons have degrees up to four and use chosen side diagonals
that the matching cap already pays, and late-pool coupons of degrees five and
six complete them (\cref{prop:record:mean-credit}).

\subsection{No positive order at the reserved cut}\label{subsec:obs:reserved}

The variance credits of \cref{sec:variance} need lower bounds for the energies
of the products of the actual half marginals of the fold records $F_1,F_2$ and
the matching records $E_0,E_1$ of a middle block. The bracket lives at
$\iota$, so the natural place to look for them is the reserved cut. It is the
wrong place.

\begin{remark}[A reserved-cut barrier and its scope]\label{rem:record:variance-n1}
At $\iota$ the half law of each of these records is folded: the six labels of a
layer take three distinct values, in the pattern of the fold
(\cref{lem:record:fold-identification}), while the root of the reference
labels a layer injectively. So the half laws are singular to that of the root,
and an independent analysis (statement N1 of
\path{research/roadmap_2026_10_18_claude/item03/main/VARIANCE_NOTE.md}) drew
three consequences: (i) no positive semidefinite order at $\iota$ between a
record's product of marginals and the reference bounds its unselected energy
from below, since in the block coordinate that product has rank one; (ii) an
exact toy with a folded law singular to the root law (energies $1$ and $1/k$ in the selected and unselected cases of that note, $Y\approx0.998$, $k=4,\dots,256$) violates the
inequality such an order would give; (iii) measuring the zero-marginal record
differences by the reference's centred norm in $L^2(m_\Blk)$ gives norms larger by a factor of about $10^3$ than the centred norms in the records' own marginals. These statements concern orders at $\iota$ against the reference,
and they are correct there. They explain how the restoration of
\cref{sec:variance} is built: it compares the root with $F_j$ at the fold
chamber of $F_j$, where both half laws live on six distinct labels with
mutually bounded densities (\cref{hyp:variance:A}, with a pointwise and a
chi-square bound in \cref{lem:record:root-marginal}), reaches $E_0,E_1$ by a
polarization at a chamber cut, and at $\iota$ uses only each record's own
centred norm in its own actual marginal space (\cref{subsec:variance:centred}).
The toy has no chamber cut and violates hypothesis (A) of \cref{hyp:variance:A} (mutually bounded half-law densities on the common chamber), which step~(B) of \cref{lem:record:restoration} uses, so it is not a counterexample to it; the
referees of the record gate reached the same conclusion.
\end{remark}

\subsection{Pairing and signed blocks: norms against envelopes}\label{subsec:obs:pairing}

The fold records $F_1,F_2$ must be replaced by laws that accept the queries and
keep the actual half marginals. The final construction uses the signed
whole-block law (\cref{def:record:mixture}): its first-order change at $\iota$
is purely skew, a three-arm star (a star with three cross pairs) of norm about $1.2732\,r$, and its whole
symmetric part is of order $1/p$ (\cref{prop:signed-bounds}). At the
October~18 checkpoint a different choice won, the orientation mixture of the
next remark, and at row~11 the pure pairing, which the signed law contains as
its case $\tlaw=0$ (\cref{lem:mirror}).

\begin{remark}[Pairing and its barrier]\label{rem:record:pairing-barrier}
The pure paired repair has leading constants $(0,\sqrt2)\,r$ (symmetric, skew)
for $F$ and $(0,1)\,r$ for the paired matching record. For $p\equiv3\pmod4$,
$p\ge43$, every accepting replacement of the two-layer $F_1$ record with the
exact old half marginals that keeps the old cross edge $y_0\to x_2$ differs from
the reference record by at least $(2-3q)^{1/2}r$ in operator norm on $L^2(m_0)$ ($q=1/p$; $m_0$ as in \cref{lem:old-f-law}),
and every accepting replacement of $E_0$ or $E_1$ by at least $r/(1-q)$ (Lemma~LB
of \path{research/roadmap_2026_10_18_claude/item01/deepen2/BARRIER.md}); so
pairing is optimal only within the class that keeps $y_0\to x_2$. The orientation mixture
of the October~18 checkpoint leaves that class: it mixes the four laws that keep
or reverse $y_0\to x_2$ and $x_0\to y_1$, with weights
$(.5803723,.2409463,.1786814,0)$ on (keep, keep), (keep, reverse),
(reverse, keep), (reverse, reverse), and its leading coefficients have norms
$0.30636\ldots$ and $0.97207\ldots$, of sum $1.2784\ldots<\sqrt2$. Joint flips
of the two edges beat pairing at first order for every ratio of symmetric to
skew price below $\sqrt{7/3}$
(\path{research/roadmap_2026_10_18/coordinator/ORIENTATION_MIXTURE_PARETO.md}).
The signed whole-block law leaves that class as well: in all three of its
components the factor $1+\mathsf B$ of the old cross edge $y_0\to x_2$ is replaced by $1+(1-\tlaw)\mathsf B$, two of them reverse the mirror edge $x_0\to y_1$, and
one of these also the old edge $x_0\to y_2$
(\cref{lem:signed-components,def:record:mixture}; master \S43,
\path{paper/sections/advance_2026_10_24/local.tex}:19-36, 95-108).
\end{remark}

At row~11 the construction returned to pure pairing, which loses on norms,
because the closing prices the symmetric and the skew parts differently. In the
reserved cost $B(e^S\cosh A-1)-D$ the symmetric price enters at first order,
through $e^S$, the skew price only through $\cosh A$ (\cref{rem:middle-bracket});
and the even skew products are charged by positive joint envelopes below $B$
(\cref{thm:envelopes,lem:reflection-averaging}), the aligned classes keeping a
fraction of their variance credit (\cref{lem:class-credit}). Against row~10,
which kept the mixture, $S$ fell from about $0.12$ to below $0.0181$ and $A$
rose from about $0.70$ to below $0.952$, and the discount $D>0.428$ more than
paid for the larger $A$ (old \S41). Two repairs are compared by their prices in
the complete reserved cost, not by the norms of their leading terms. The signed
law of the final construction keeps the zero symmetric leading change of
pairing and lowers its skew leading change from $\sqrt2\,r$ to about
$1.2732\,r$, at the variation $1+\tlaw$ (\cref{rem:obs:sign-frontier}); the pure
law itself had reached its limits (\cref{rem:obs:pure-sharpness}).

\begin{remark}[The pure paired law at its limits]\label{rem:obs:pure-sharpness}
Take the pure paired law in the form of row~12: its raw kernel is that of the
case $\tlaw=0$ of \cref{def:record:mixture}, and its repair and its actual
leading two-arm star $L_{F,p}$ are those of master \S42. Its difference at the prime $p$, $D_{F,p}$, satisfies, for primes $p\equiv3\pmod4$,
\[
 \lim_{p\to\infty}p\,\|D^{\rm s}_{F,p}\|=2,\qquad
 \lim_{p\to\infty}p\,\|D^{\rm a}_{F,p}-L_{F,p}\|=\sqrt5,
\]
with witnesses on nontrivial translation sectors (master \S42,
\texttt{cont:local:sharpness},
\path{paper/sections/continuation_2026_10_24/local.tex}:1074-1156;
\path{research/continuation_2026_10_24/local/SHARPNESS.md}); at the first piece
of row~12 the two coefficients were about $2.010$ and $2.268$, against about
$10.836$ and $16.677$ in the bounds of row~11. The coefficient $\sqrt6$ of the
direct old/new Gram of the matching envelopes of that law is sharp for scalar
side diagonals: the leading two-side Gram has the characteristic polynomial
$z^4(z^2-6)(z^2-2)$, so a common scalar side shift must be at least $\sqrt6$,
and two different scalar shifts have average at least $\sqrt6$
(\path{research/continuation_2026_10_24/local/DIRECT_GRAM_OBSTRUCTION.md}).
Both statements concern that law and that decomposition: a further gain in the
middle needed another law, and they say nothing about the optimality of the
exponent.
\end{remark}

\begin{remark}[The leading frontier of signed fold laws]\label{rem:obs:sign-frontier}
Average the $256$ complete cross-sign patterns of the fold record that contain
the accepted query, with arbitrary real weights of total mass one, and let
$n_{\rm s}$, $n_{\rm a}$ be the symmetric and skew leading norms of the change, in
units of $r$. Then, for every relative price weight $c\ge1$ on the symmetric part, $c\,n_{\rm s}+n_{\rm a}\ge\sqrt{3/2}$, without
any positivity. A zero symmetric leading error forces the four sign moments of
the cross edges into a family with one parameter $\tlaw$, whose skew leading
norm is $(1+\tlaw^2+(1-\tlaw)^2)^{1/2}$, least at $\tlaw=1/2$; nonnegative weights
force $\tlaw=0$, pure pairing, with $\sqrt2$; for $0\le\tlaw\le1$ a signed law of the
family has total variation at least $1+\tlaw$, and the three-component law of
\cref{def:record:mixture} attains it. On the same family the scalar coefficient of the direct Gram of the matching sides in the leading Euclidean model is at least
$s_\tlaw=(4-2\tlaw+2\tlaw^2+2(1+\tlaw^2)^{1/2})^{1/2}$, with equality for that law
(\cref{lem:record:direct-gram}): $\sqrt6$ at $\tlaw=0$, about $2.3842$ at
$\tlaw=127/500$
(\path{research/advance_2026_10_24/local/SIGNED_LEADING_FRONTIER.md}:9-82, checked
by \path{research/advance_2026_10_24/local/check_signed_frontier.json}; master
\S43, \texttt{adv:local:signed-law},
\path{paper/sections/advance_2026_10_24/local.tex}:95-146). The statement
concerns the leading Euclidean model: it rules out a further leading gain from
more signed cross orientations with these objectives, not a different law, and
the selected $\tlaw=127/500$ is chosen by the complete closing, which also depends
on the envelopes, the restorations, the coupons and the finite remainders.
\end{remark}

\subsection{The walls of the prefix at primes \texorpdfstring{$p\equiv1\pmod4$}{p = 1 (mod 4)}}\label{subsec:obs:prefix}

A \emph{wall} in this section means a value of the exponent at
which one explicit bound of the construction reaches $1$, so that
this layout of the construction cannot pass it without a new idea.
For the prefix the binding bound is the price of deleting
coincident labels from the square clique; the remark computes the
two walls it gives.

\begin{remark}[The prefix walls]\label{rem:record:prefix-walls}
The numerical comparisons in this remark use the baseline $\vartheta=1$ and $a=\PushA$. With $\mcl$ digit pairs the prefix deletes its exception at a total price at most
$\eta_{\mcl}P_4^{1+2\mcl a}$, $\eta_{\mcl}=2^{1-\mcl}-2^{-2\mcl}$
(\cref{lem:frontier:iidenergy}),
which at $P_4=2^{512}$ is below $2^{513-\mcl(1-1024a)}$: each digit pair halves the
exception and multiplies the selection factor of a prefix prime by at most
$2^{1024a}$. This gives two walls. (i) For a free number of digit pairs the
bound tends to zero exactly when $2^{1024a}<2$, that is $a<1/1024$, or
$\eps<1/72704\approx1.3754\cdot10^{-5}$; at the baseline exponent
$1024a=0.909759\ldots$ and $2^{1024a}\approx1.879$, so $\eps$ is about $91\%$ of
this wall. (ii) With $\mcl=2048$, the number used up to row~10, the exponent
vanishes at $a=1535/2^{21}$, that is $\eps\approx1.0309\cdot10^{-5}$. The October~18 checkpoint (binary exponent $513-\mcl(1-1024a)$ about $-55.2$) and row~10
($\eps\approx9.963\cdot10^{-6}$) lie below this wall; row~11 and all later rows
lie above it, and at the baseline exponent it would be about $+328.19$. This is why
$\mcl=8192$: at the baseline $a$ the exponent is negative exactly for $\mcl\ge5685$,
and at $\mcl=8192$ it is about $-226.25$, far inside the allowance $10^{-12}$;
with $8192$ pairs it stays inside the allowance on the whole analytic range
$a\le9/10000$ (at $a=9/10000$ it is $-80783/625\approx-129.25$), and it vanishes
at $a=7679/2^{23}\approx9.154\cdot10^{-4}$. More digit pairs cost only a larger
fixed modulus (\cref{lem:frontier:iidenergy}(iv), \cref{rem:prefix-digits}).
Both walls belong to the layout, not to the method:
a clique prefix stopped at a smaller power of two $P$, with the all-prime tail
starting there, replaces $1024a$ by $2a\log_2P$ and moves both walls to larger
exponents.
\end{remark}

\subsection{Why class one is a prefix}\label{subsec:obs:class-one}

The five-type completion of the middle range expresses the $F$ and $E$ kernels
against one conditional reference, the common middle law of two labels
(\cref{lem:five-type,lem:front12:reference}), and needs that law to be uniform on
ordered distinct pairs. For $p\equiv3\pmod4$ it is, because the exchange
$(u,u')\mapsto(u',u)$, which preserves the completion counts, interchanges the
two orbits of ordered distinct pairs under the maps $z\mapsto\alpha z+\beta$ with
$\chi(\alpha)=1$ (\cref{lem:record:class-three}(ii)). For $p\equiv1\pmod4$ the exchange preserves
the orbits and the counts differ: the pairs with $\chi(u'-u)=1$ have $312$ completions to $F$-configurations and the others $264$ at $p=13$ (\cref{rem:middle-class-three}), and
$1464$ against $1344$ at $p=17$ (by enumeration). At class-one primes, on the
other hand, the strict-square relation is symmetric and the square clique costs
only an exception (\cref{lem:frontier:iidlocal}). So every class-one prime below
$P_4$ is a prefix prime, the middle and upper ranges are class three
(\cref{sec:frontier:iidprefix}), and class one returns only in the tail, with
caps certified for both classes (\cref{subsec:advance:parity,sec:tail-index}).

\subsection{The root law: forced support and forced norm}\label{subsec:obs:root}

The chamber cap $d$ of the twelve-label root, the negative cap of its cut L20,
fixes the first-order weights $w_{\rm D}=w_{\rm S}$ of the reference, in
particular the weight $w_{\rm D}r$ of every fold record, through the completion
inequality $w_{\rm D}\ge d$ (\cref{lem:front12:reference}). Two scoped facts say
where a gain can come from.

\begin{remark}[Support closure and the forced norm]\label{rem:obs:root-support}
(i) Let $\Gamma$ be a probability law on the labellings of the twelve vertices
$(u,\ell)$ of $G_{\rm L,2}$ at a prime $p\equiv3\pmod4$ whose half marginal on the six half vertices $\{0,1,2\}\times\{0,1\}$ is that of the uniform root, and which is
invariant under the three chamber reflections $u\mapsto k-u$, $k\in\{1,3,5\}$.
Then $\Gamma$ is supported on the labellings that satisfy all $24$ edges of
$G_{\rm L,2}$; in particular none of the sixteen cross edges of L20 can be
dropped. Indeed the group generated by the reflections carries $\{0,1,2\}$ to the
six consecutive triples $H_a=\{a,a+1,a+2\}$ (the reflections give $H_5,H_1,H_3$,
and the translations by $0,2,4$, products of two reflections, give $H_0,H_2,H_4$);
each contains exactly one of the six rotation edges, with
its four layer lifts, and every rotation edge lies in exactly one $H_a$. The half
marginal of the uniform root is carried by the labellings of $G_{\rm F,2}$
(\cref{lem:layer12-geometry}(b)), so the failure of the edge inside $\{0,1,2\}$
is a null event, and invariance carries this to every $H_a$ (master \S43,
\texttt{adv:local:root-support-closure},
\path{paper/sections/advance_2026_10_24/local_interim.tex}:128-151). Positive
reweightings of the accepted support are allowed; signed measures and changed
half laws are not covered.

(ii) Let $T$ be the density operator, relative to $\mu_p\otimes\mu_p$ with
$\mu_p=\lambda_p(G_{\rm F,2})$, of any probability law on twelve distinct labels
that satisfies the sixteen cross constraints of L20, whatever its marginals,
and let $P_0$ be the projection onto the mean-zero functions. Then
$\|P_0TP_0\|\ge\sqrt{30}\,r\,(1-o(1))$; the proof compares normalized traces of
$T$ with the fifteen star features of \cref{lem:general-star}
(master \S43, \texttt{adv:local:forced-root-norm},
\path{paper/sections/advance_2026_10_24/local_interim.tex}:153-213). The record
\path{research/advance_2026_10_24/local/check_l20_forced_norm.json} evaluates the
finite form of the bound at four endpoints; at $Q_3$ it gives
$\|P_0TP_0\|\ge5.4472\,r$. The bound concerns the centred norm only: a diagonal entry of order $r$ not constrained by the cross conditions in a higher-star block can move the negative
eigenvalue of that block without changing its forced off-diagonal entries, so the
negative eigenvalue is not forced.
\end{remark}

So a law that keeps the symmetries and the half marginal of the root cannot drop
a constraint, and a gain must come from a positive reweighting of the accepted
support; by (ii) it can lower only the negative cap, the centred norm of every
law on that support being at least about $\sqrt{30}\,r$. The completion
inequalities and the block caps use the negative
and the positive cap at each cut (\cref{lem:layer12-caps}), not the centred
norm, and a reweighting that changes the half marginal of the root is paid in
the restorations (\cref{lem:record:root-marginal}(iv)). This is what the
positive roots of \cref{sec:positive-roots} do. In the selected construction,
the chamber block caps $w_{\rm D}$ on the pieces from $2^{36}$, $2^{40}$ and
$2^{45}$ are below $5.46078$, $5.45753$ and $5.43446$, against about
$5.4776$, $5.4773$ and $5.4772$ for the uniform root
(\cref{cert:positive-root-caps,cert:l20-grid}); four earlier pieces are tilted
as well. Their centred norm still satisfies (ii), so the gain concerns the
negative cap that the completion actually pays. The asymptotic tilts
that preceded them (an eighth-trace tilt whose centred quartic diagonal response
is at least $0.9065\,r$ for $p\ge2^{20}$; quartet tilts with limiting
coefficients about $5.4718$ and $5.4754$; a linear pair-trace tilt, about
$5.4664$; the average of individual squares, about $5.4562$; the square of the
average, about $5.4421$) carried unpriced $O(q)$ remainders and bounds on their
own marginals, and none moved the exponent until both cuts and the restoration
inputs were proved in the fixed charts (master \S43,
\path{paper/sections/advance_2026_10_24/local_interim.tex}:215-1349;
\cref{subsec:history:oct2425}).

\subsection{Why twelve queries do not help}\label{subsec:obs:twelve}

The twelve queries of $\mathcal E_{12}$, with the sources $(u,v,0)$,
$u\in\{0,5\}$, have $24$ distinct endpoints and project to two first-coordinate
pairs (\cref{subsec:front12:construction}). With twelve disjoint indices each
conductor still occurs twice in the complementary folds, the least square step
is unchanged, and the clique exception is the two-pair mass already charged in \cref{lem:frontier:iidenergy}. So,
given comparison data for $\mathcal E_{12}$ with the actual marginals and the
bounds of (CP0)--(CP8), \cref{thm:comparison} with $k=12$ applies, and its closing
condition reads
\[
 \frac{4095}{4096}\bigl(L_0+\mathcal E\bigr)+J_{12}+W_{12}+\frac{1-m}{4096}<1,
 \qquad J_{12}=Be^{T_{12}}\Bigl(\frac{R_I}2\Bigr)^{12},
\]
\[
 W_{12}=Be^{S+A}\Bigl[\Bigl(1+\frac{R_I}{2\Kint^{35/36}}\Bigr)^{12}-1\Bigr],
\]
with $T_{12}$ the matching exponent of such data (master \S43,
\texttt{adv:upper:twelve-interface},
\path{paper/sections/advance_2026_10_24/upper_twelve.tex}:27-96). Such data are not
available. Locally, accepting laws exist: exact Markov couplings of the actual
half law accept both fold queries $x_1\to y_0$, $x_2\to y_1$, or both queries of
a matching fork, and their operators vanish on invariant mean-zero functions, so
\cref{cor:markov-completion} with such a coupling turns every nonnegative
invariant accepting raw kernel of mass one with bounded marginal defects ($|\alpha|,|\beta|\le M$ there) into a
law with exactly the actual half marginals (master \S43,
\texttt{adv:local:two-query-markov}, \texttt{adv:local:two-query-repair},
\path{paper/sections/advance_2026_10_24/local_interim.tex}:1-120). The
obstruction is the price on the unchanged references.

\begin{remark}[Two queries force a symmetric error]\label{rem:obs:twelve}
Let $K$ be any real signed law of mass one on pairs of halves of the fold record,
written as a density relative to $m_0\otimes m_0$, with $m_0$ the actual half
marginal of \cref{lem:old-f-law}, that accepts $\chi(y_0-x_1)=\chi(y_1-x_2)=1$.
Then $\|P_0(\operatorname{Sym}K-T_{\rm old})P_0\|\ge(1/\sqrt3-o(1))\,r$, and uniformly
at least $0.57263\,r$ for $p\ge2^{24}$, where $T_{\rm old}$ is the reference fold
record of \cref{lem:old-f-law} and $P_0$ the projection off the constants of
$L^2(m_0)$. For the two queries with common source $x_0\to y_1$, $x_0\to y_2$ (the matching fork), every signed law of mass
one on uniform injective triples that accepts $x_0\to y_1$ and $x_0\to y_2$ has
symmetric and skew changes whose norms add up to at least $(\sqrt3-o(1))\,r$
(master \S43, \texttt{adv:upper:actual-two-query},
\texttt{adv:upper:actual-fork-frontier},
\path{paper/sections/advance_2026_10_24/upper_interim.tex}:209-370;
\path{research/advance_2026_10_24/upper/check_actual_two_query_barrier.json}). No
positivity, invariance or marginal condition is assumed. With these references a
twelve-query fold record pays a symmetric price of order $r$, not $1/p$, and a
twelve-query matching record with a symmetric error of order $1/p$ has a skew
coefficient at least about $\sqrt3$. The statement concerns the unchanged half
laws and reference records, not twelve queries as such: a different reference,
or a joint treatment of the symmetric terms, could change its consequences.
\end{remark}

A twelve-query construction would therefore need fold and matching records that
accept both queries of their halves, with comparison bounds against a changed
reference, and new class bounds (CP5) for their leading stars. Parts of the
construction accept $\mathcal E_{12}$ already, but its six-query refinements do
not: the signed fold records of the middle accept, in each fold, only the one
cross edge that $\mathcal E_6$ needs there, while $\mathcal E_{12}$ needs both
$x_1\to y_0$ and $x_2\to y_1$; the LE components of the upper range add only the
edge that $\mathcal E_6$ requires (\cref{def:dense-comparison,prop:le-mixture}); and the
records $G_{01}L$, $G_{10}L$ are retained because they accept $\mathcal E_6$
(\cref{lem:six-query-retention};
\path{research/advance_2026_10_24/upper/SIX_QUERY_UPPER.md}). A twelve-query construction would have to repair or delete them instead of retaining them.

\subsection{Obstruction polynomials of the singleton recurrence}
\label{subsec:obs:polynomials}

Before the comparison theorem, records that did not accept a coupled family
were covered, not repaired. In the singleton-recurrence design of the
$216$-slot period (\cref{subsec:history:216}), nine edge families crossing the reserved cut are
indexed by $\Phi=\{0,1,2\}^2$, with rows $\mathcal R_i=\{(i,j):j\}$ and columns
$\mathcal C_j=\{(i,j):i\}$. The law of a middle atom is a sum of typed parts:
type $\star$ accepts every family, a residual part none, and the fifteen
\emph{rescue types} $\mathrm{row}_i$, $\mathrm{col}_j$, $\mathrm{cr}_{i,j}$
accept $\operatorname{Acc}(\mathrm{row}_i)=\mathcal R_i$,
$\operatorname{Acc}(\mathrm{col}_j)=\mathcal C_j$ and
$\operatorname{Acc}(\mathrm{cr}_{i,j})=\mathcal R_i\cup\mathcal C_j$. A record,
one type at every atom, accepts the intersection of the acceptance sets of its
types. For a coupled set $\emptyset\ne\mathfrak F\subseteq\Phi$, records that
accept a family of $\mathfrak F$ cancel through the actual-square kernel; the
others were covered in the Loewner order at the reserved cut, which fixes every
middle record samplewise, by one \emph{cylinder} per residual
part and one per minimal obstruction placed on distinct atoms,
where a cylinder here is a law depending only on the labels at
finitely many named atoms (a cylinder set in the product). A set $S$ of rescue types is
\emph{$\mathfrak F$-incompatible} if
$\mathfrak F\cap\bigcap_{t\in S}\operatorname{Acc}(t)=\emptyset$ (the empty
intersection being $\Phi$), and a \emph{minimal $\mathfrak F$-obstruction} is an
inclusion-minimal incompatible set. With priced intensities at most $x$, $y$,
$z$ for rows, columns and crosses, the cover costs at most
$Z+P_{\mathfrak F}(x,y,z)$, with $Z$ the residual intensity and
$P_{\mathfrak F}=\sum_Sx^{n_{\rm r}(S)}y^{n_{\rm c}(S)}z^{n_{\rm x}(S)}$ over the
minimal obstructions, $n_{\rm r},n_{\rm c},n_{\rm x}$ counting the rows, columns
and crosses of $S$. The all-family criterion closed when
$\frac12M^{|\mathfrak F|}[Z+P_{\mathfrak F}+|\mathfrak F|/\Kint+e^H]<1$, with $M$
the modulus moment of the index and $H$ the upper exceptional intensity. The
singleton checkpoint ($\eps\approx3.3\cdot10^{-9}$) coupled the nine families
one at a time and covered the records that accept none of them, with
$P_\Phi$ about $0.369$.

\begin{lemma}[Rejection-class covers]\label{lem:obs:rejection-covers}
Let $\emptyset\ne\mathfrak F\subseteq\Phi$, put
$\operatorname{Rej}(t)=\mathfrak F\setminus\operatorname{Acc}(t)$ for a rescue
type $t$, give $\mathrm{row}_i$, $\mathrm{col}_j$, $\mathrm{cr}_{i,j}$ the weights
$x,y,z$, and let $w_X$ be the total weight of the types $t$ with
$\operatorname{Rej}(t)=X$. A \emph{minimal cover} of $\mathfrak F$ is a family of
distinct sets with union $\mathfrak F$ no proper subfamily of which has union
$\mathfrak F$.
\begin{enumerate}[label=(\roman*)]
\item $S$ is $\mathfrak F$-incompatible if and only if
 $\bigcup_{t\in S}\operatorname{Rej}(t)=\mathfrak F$.
\item $S$ is a minimal $\mathfrak F$-obstruction if and only if the sets
 $\operatorname{Rej}(t)$, $t\in S$, are nonempty, pairwise distinct, and form a
 minimal cover of $\mathfrak F$.
\item $P_{\mathfrak F}=\sum_{\mathcal C}\prod_{X\in\mathcal C}w_X$ over the minimal
 covers $\mathcal C$ of $\mathfrak F$ by nonempty rejection sets, and every
 minimal $\mathfrak F$-obstruction has at most $|\mathfrak F|$ elements.
\end{enumerate}
\end{lemma}

\begin{proof}
(i) is De Morgan's law,
$\mathfrak F\cap\bigcap_S\operatorname{Acc}(t)=\mathfrak F\setminus\bigcup_S
\operatorname{Rej}(t)$; for $S=\emptyset$ both sides are false. (ii) By (i)
compatibility passes to subsets, so $S$ is minimal if and only if it is
incompatible and every $S\setminus\{t\}$ is compatible. If $S$ is minimal, no
$t\in S$ has $\operatorname{Rej}(t)=\emptyset$ and no two share a rejection set,
since removing such a $t$ would not change the union; a proper subfamily of
$\{\operatorname{Rej}(t):t\in S\}$ is then the family of a proper subset of $S$,
which does not cover. Conversely such a family makes $S$ incompatible and every
$S\setminus\{t\}$ compatible. (iii) By (ii) the minimal obstructions correspond
bijectively to the pairs of a minimal cover $\mathcal C$ of $\mathfrak F$ by
nonempty rejection sets and a choice of one type $t_X$ with
$\operatorname{Rej}(t_X)=X$ for each $X\in\mathcal C$; for fixed $\mathcal C$,
summing the monomials over the choices gives $\prod_{X\in\mathcal C}w_X$. Each member of a minimal cover
has a point of $\mathfrak F$ in no other member, so $|\mathcal C|\le|\mathfrak F|$.
\end{proof}

For the transversal $\mathfrak D=\{(0,0),(1,1),(2,2)\}$ the types of
$A_i=\{\mathrm{row}_i,\mathrm{col}_i,\mathrm{cr}_{i,i}\}$ reject the two
diagonal cells other than $(i,i)$, the two crosses of
$B_k=\{\mathrm{cr}_{a,b},\mathrm{cr}_{b,a}\}$, $\{a,b,k\}=\{0,1,2\}$, reject
$(k,k)$; with the weights $s=x+y+z$ and $2z$ and the minimal covers
$\{A_i,A_j\}$, $\{A_i,B_i\}$ (three each) and $\{B_0,B_1,B_2\}$,
$P_{\mathfrak D}=3s^2+6sz+8z^3$. For the two-pair set
$\mathfrak E=\{(0,1),(1,0)\}$ the types $\mathrm{row}_2,\mathrm{col}_2,
\mathrm{cr}_{2,2}$ reject $\mathfrak E$, two classes of weight $x+y+3z$ reject
one cell each, and two crosses reject nothing: $P_{\mathfrak E}=s+(x+y+3z)^2$.

\begin{table}[ht]
\centering
\small
\caption{Obstruction polynomials of four coupled sets, $s=x+y+z$
(\cref{lem:obs:rejection-covers}; confirmed by enumerating all sets of rescue
types).}
\label{tab:obs:polynomials}
\begin{tabular}{@{}>{\raggedright\arraybackslash}p{0.2\textwidth}ccc
                 >{\raggedright\arraybackslash}p{0.47\textwidth}@{}}
\toprule
Coupled set $\mathfrak F$ & $|\mathfrak F|$ & minimal & largest &
$P_{\mathfrak F}(x,y,z)$\\
\midrule
$\Phi$ (all families) & $9$ & $165$ & $4$ &
$3x^2+3y^2+36xyz+36(x+y)z^2+6z^3+45z^4$\\
transversal $\mathfrak D$ & $3$ & $53$ & $3$ & $3s^2+6sz+8z^3$\\
two-pair set $\mathfrak E$ & $2$ & $28$ & $2$ & $s+(x+y+3z)^2$\\
two rows $\mathfrak R=\mathcal R_0\cup\mathcal R_1$ & $6$ & $77$ & $4$ &
$x+x^2+3(y+z)^2+12xyz+18xz^2$\newline${}+12(y+z)z^2+9z^4$\\
\bottomrule
\end{tabular}
\end{table}

The two-column set is the transpose of $\mathfrak R$, with polynomial
$P_{\mathfrak R}(y,x,z)$. The transversal and the two-pair set occurred only in
research designs, the two-row set in an earlier unpublished design. By
\cref{lem:obs:rejection-covers}(i), $P_{\mathfrak F}$ has a term of degree one
exactly when $\mathfrak F$ misses a row or a column (if
$\mathcal R_i\cup\mathcal C_j$ misses $\mathfrak F$, so does $\mathcal R_i$). For
one family, $\mathfrak F=\{(i,j)\}$, the two other rows, the two other columns
and the four crosses off row $i$ and column $j$ are obstructions by themselves:
$P_{\{(i,j)\}}=2x+2y+4z$. This first-order term is why coupling a single pair is
costly; at the type masses of the first version's middle terminal it was about
$61.6$ per unit of terminal weight. A set without first-order terms has at least
three cells, with equality exactly for the six permutation matrices, which all
have $P_{\mathfrak D}$ by the symmetry $S_3\times S_3$ of $\Phi$; but the
modulus moment enters as $M^{|\mathfrak F|}$, so larger sets cost more in the
index. Which set paid depended on the design. Theorem~C has no obstruction
polynomial: its queries are fixed, and a record that does not accept them is
repaired, keeping the actual conditional marginals, or deleted at the price of
its own mass (\cref{sec:completion-repair}). A record of weight $O(r)$ whose
repair moves it by $O(r)$ then costs $O(1/p)$ instead of entering a union bound
with its full intensity; this was the source of the gain of the fixed-family
transfer, the $216$-slot construction that is the first instance of the scheme
of Theorem~C (\cref{subsec:history:216}).

\subsection{Limits of the method}\label{subsec:obs:limits}

\emph{The kernel cutoff.} The circle method of \cref{app:kernel20} needs its
major-arc exponent in the window $(10\beta/3,\,1/4-19\beta/12)$, nonempty exactly
when $\beta<3/59$ (\cref{rem:record:K20-range}); at $\beta=1/20$ the window is
$(1/6,41/240)$ and contains $17/100$. So moduli up to $N^{1/20}$ are close to the
limit $N^{3/59}$ of this argument.

\emph{The error exponent.} For the baseline, the induction needs $a<\sigma$
(\cref{thm:record:lineardensity}), so the stage of
\cref{def:record:analytic-stage}, $\sigma=1/825$, cannot give more than
$\eps<1/(71\cdot825)\approx1.707\cdot10^{-5}$, whatever the local laws. It is not
binding: $\sigma-a=65195254/201416015625$, $a\approx0.733\,\sigma$, and the
exponent is limited by $\Theta<1$. A larger $\sigma$ needs a new stage, since $\sigma$ enters (S1)--(S4) of the stage inequalities (\cref{def:record:analytic-stage}), with new kernel and analytic
certificates; this was not attempted. The gate $a<\sigma$ does not hinge on the
kernel at $\beta=1/20$: with a separate outside cutoff the comparison needs no
kernel and the cancellation needs it only at $\beta=1/36$, and the stage
$(\beta,\tau,t_{\rm an},\sigma)=(1/36,1/230,1/108,1/1000)$ was checked at the
October~18 exponent; it is not used.

\emph{The analytic domain.} For the baseline, the analytic stage, the transfer gates and the
prefix allowance are certified on the whole range $0<a\le9/10000$, independently
of $Q_3$ and $P_3$ (\cref{cert:record:analytic-stage};
\path{research/advance_2026_10_24/global/ANALYTIC_DOMAIN.md}). This allows
$\eps\le9/710000\approx1.2676\cdot10^{-5}$, about $1.30\%$ above the baseline
exponent; with $8192$ digit pairs the prefix wall lies at
$\eps\approx1.2893\cdot10^{-5}$ (\cref{rem:record:prefix-walls}). Beyond the range
the gates must be audited again, and beyond that wall the prefix needs more
digit pairs; this is a limit of the certificate and of the layout, not of the
method, and the baseline closing becomes limiting before it is reached
(\path{research/further_2026_10_25/SYNTHESIS.md}:187-195). The selected theorem rechecks the analytic, prefix and density gates on $0<a\le91/100000$ and uses $\vartheta=\FourRoutesDensity$ (\cref{cert:four-routes-closing}).

\emph{The $2$-adic scope.} The earlier unit-square kernel used a $2$-part
of $8$. The root classes of \cref{lem:kernel-depth-classes} now cover every
nonvanishing square class with at least three remaining binary digits after
its even valuation is removed. This includes the selected physical modulus
$2^{129}$ and root valuations below $64$. Zero classes are still excluded;
conditioning on the random depth does not replace the full physical modulus $\Mzero$ in cancellation (\cref{rem:record:item13}).

\emph{Second-order inputs.} The certified budget uses Deligne's theorems, the
Katz--Sarnak local monodromy and the evaluation of Ahlgren and Ono
(\cref{thm:ahlgren-ono}) for its second-order bounds. This explains their role in
this proof; it does not establish that they are necessary for a comparable
exponent by another argument (\cref{subsec:intro:deligne}).

\emph{Allowed couplings in the envelope bounds.} One natural off-diagonal operator
between fold and matching features, a coherent cross arm, has a kernel whose
total variation, and hence the coupling mass of every expansion of it into
positive couplings, is at least $(p-2-1/p)/(2\sqrt p)$, unbounded in $p$; so it
cannot enter the uniform comparison of the envelopes at no extra price
(\path{research/advance_2026_10_24/envelopes/COHERENT_ARM_OBSTRUCTION.md}). The
larger invariant background of the bracket with one fold copy improves the
exponent of the invariant envelope cap but supplies no new permutation of the
fold features, since every element of the root group maps a fold pattern to a
fold pattern with the identity permutation of its object names
(\path{research/advance_2026_10_24/envelopes/PHYSICAL_ORBIT_SCOPE.md}). The
matching energy pairs a test function on the source side with one on the target side without swapping them, so the envelope reflection is not available there for free
(\path{research/advance_2026_10_24/envelopes/MATCHING_REFLECTION_SCOPE.md}). And in
the pure-$E$ singleton algebra every positive $Z\succeq\pm T$, $T$ a tensor of $k$
unit matching arms, has collected coupling mass at least $4^k$, so the uniform
discount of a pure-$E$ class cannot be positive once
$(R_{\rm err}-2u)4^k\ge B(A_E/A^{\rm arm}_E)^k$, $u$ the exactly-one-tag credit per
unit mass (master \S43, \texttt{adv:env:pure-E-mass-floor},
\path{paper/sections/advance_2026_10_24/envelopes_interim.tex}:214-249). These are
limits of the present allowed couplings, not of positive envelopes.

\emph{Fold copies in the bracket.} The bracket carries one copy of the folded
fold records (\cref{lem:record:gram}). With $c<1$ copies the protection
inequality, which then reads $c\,v_{\rm D}\ge2k_{\rm D}w_{\rm D}$ (it is
\eqref{eq:protection} at $c=1$), raises $v_{\rm D}$; the reference with $c<1$ copies is, before renormalization, a subprobability (total mass $<1$) whose normalization enters every comparison coefficient; and
subtracting only the two selected fold patterns breaks the root invariance on
which the invariant envelope caps rest. These are requirements of a different
reference, not an impossibility
(\path{research/further_2026_10_25/envelopes/BETA_BELOW_ONE_SCOPE.md}).

\emph{The layer cap of the positive roots.} For the individual-square family,
the proved layer coefficient tends, as $p\to\infty$, to $\sqrt3+\sqrt7$ for each fixed $\ulaw>0$, twice the norm
$(\sqrt3+\sqrt7)/2\approx2.189$ of the leading matrix of one component of the
layer cut; the layer caps of the uniform root on these pieces are between about
$2.189$ and $2.196$ (\cref{cert:layer12-grid}). The proof prices the two
one-component terms of the tilted kernel separately, through the auxiliary
space, so its bound does not tend to the uniform value as $\ulaw\to0$
(\cref{prop:root-layer};
\path{research/further_2026_10_25/local/L21_PAIR_SQUARE_TILT.md}:84-125, 240-250).
The selected estimates bound the finite-$p$ remainder of this argument directly. A sharper
leading layer estimate for this family would lower the weight $v_{\rm L}$
on its six tilted pieces; the coherent last-piece law has its own layer
bound in \cref{sec:positive-roots}.

\emph{Six queries.} With the pure paired comparison law six disjoint universal
queries are the most: a universally supported query must project to one of the
six directed rotation edges of the root, and on the common support of the root and of the three fold records $F_0,F_1,F_2$ the only universally supported queries run from
layer $0$ to layer $1$ over the first-coordinate pair $0\to2$, so each of them
meets one of the six slots $(0,v,0)$
(\path{research/continuation_2026_10_24/envelopes/QUERY_OBSTRUCTION.md}). This is
a support statement for that law, not for every modified comparison law;
\cref{subsec:obs:twelve} says why the twelve-query family is not used either.

\subsection{Tested and not pursued}\label{subsec:obs:tested}

\Cref{sec:history} places these in the development; here is why each was
dropped.
\begin{itemize}
\item \emph{An E-only envelope of degree six at row~11.} Its coupling mass was
 about $4103.8$, against about $257.9$ with four factors, and at the selected
 endpoint of row~11 its complete cost exceeded $B_*$, so the coupling error
 consumed the saving. In the final portfolio the class $(0,6)$ pays, with a
 saving of about $4.9\cdot10^{-6}$ (\cref{cor:envelope-costs}); the coupling-mass
 floor of the pure-$E$ classes is in \cref{subsec:obs:limits}.
\item \emph{Graph laws at class-one primes.} The octahedron law on $\Zsix$
 (twelve unordered strict-square edges, all pairs but the antipodal ones, which
 make sense because $-1$ is a square) has small completion weights, but query
 edges on which its records have zero difference needed fine indices and their
 energy factor; a larger law with strictly square accepted edges served the upper
 part of the finite prefix. The square clique costs only an exception.
\item \emph{The sparse middle law and the best-cut comparison} of the
 $216$-slot period ($180$ of the $252$ middle edges, one amplified edge and then
 four, every term estimated at its best cut; \cref{subsec:history:sparse}). In
 the final design every query
 crosses the reserved cut; an ordinary cut and the reserved cut need not share a
 comparison energy, and a compatible multi-cut theorem was not proved.
\item \emph{An averaged completion denominator for the upper root} gives $g\approx4887$ for the root cut LL20 at $2^{49}$, computed from inputs that were not certified, against $2711.73\ldots$ for the direct-star root
 (\cref{thm:record:direct-star}) at the October~18 checkpoint; in this
 architecture the degree-twenty root is forced.
\item \emph{The other directions of the October~23 handoffs}
 (\path{research/push_2026_10_24/README.md}): stronger upper reports that could
 not be recovered and changed extraction and density frameworks; rows 12--14
 took up further background and mean savings in their own form
 (\cref{subsec:history:oct2425}). Their gains do not add to the final result
 without compatible laws and complete costs.
\item \emph{A continuous orientation family} of the pure law at row~12, mixing
 its four orientation components, each repaired to the same actual marginal:
 every direction on the finite grid of mixture weights tested in the two named records worsened the complete bound, and the tests are not an optimality proof
 (\path{research/continuation_2026_10_24/local/ORIENTATION_CONTINUATION.md},
 \path{research/continuation_2026_10_24/README.md}:128-135).
\item \emph{Quartet tilts of the root.} The one-cut and the chamber-average
 quartic tilts have limiting coefficients about $5.4718$ and $5.4754$, above
 those of the pair-trace tilts (master \S43,
 \texttt{adv:local:single-cut-negative-tilt},
 \texttt{adv:local:chamber-negative-tilt},
 \path{paper/sections/advance_2026_10_24/local_interim.tex}:487-881).
\item \emph{Larger kernels and pinned Grams.} The complete four-vertex character
 kernel on injective quadruples has norm at most $24p^{-2}$ (previously
 $96p^{-2}$; master \S43, \texttt{adv:upper:quartic24},
 \path{paper/sections/advance_2026_10_24/upper_interim.tex}:9-179), but the selected extraction class counts do not change, and the bound does not justify
 deleting collision terms of the grouped extraction (\cref{subsec:obs:pitfalls}).
 One-pin, same-pin and two-pin weighted pair Grams through a twisted
 symmetric-square Legendre moment (\texttt{adv:env:one-pin-pair},
 \texttt{adv:env:same-pin-conditional}, \texttt{adv:env:two-pin-pair},
 \path{paper/sections/advance_2026_10_24/envelopes_interim.tex}:9-173) give a
 conditional pair-arm corollary that no selected row uses, since none has that
 pin structure. The absolute two-tag composition
 (\texttt{adv:env:absolute-two-tag}, envelopes\_interim.tex:175-212) relates its
 endpoint energies by a map that is not a chamber reflection of the reference,
 so the grouped first-stage error keeps its countwise minimum.
\item \emph{Proved refinements not used in the closing.} Support-dependent conditional
 densities improve the conditional $F$ row constants, but other completion
 constraints already determine the selected weights (master \S43,
 \texttt{adv:local:conditional-support-density},
 \path{paper/sections/advance_2026_10_24/local.tex}:1594-1676); the even-degree
 root-range parameters $10r+7q$, $10r+71q$, $10r+103q$, $10r+31q$, which
 \cref{thm:weil}(b) gives, are used by no scalar (\texttt{adv:global:root-even},
 \path{paper/sections/advance_2026_10_24/global.tex}:178-199); the one-edge and
 two-edge squared-trace kernels with leading constant $4$ (master \S44,
 \texttt{fur:local:one-edge}, \texttt{fur:local:two-edge},
 \path{paper/sections/further_2026_10_25/local.tex}:1-559) aim at a balanced root
 law of $48$ scores whose cross-score covariance is unproved.
\item \emph{Other root laws.} A quartic multiplier model suggests a coefficient
 about $5.4229$, conditional on unproved cubic and quartic identities for stars rooted at one vertex
 (\path{research/further_2026_10_25/global/QUARTIC_TILT_MODEL.md}); a balanced
 two-orbit model of $48$ scores about $5.3961$, with its rooted covariance
 unproved (\path{research/advance_2026_10_24/WORK_LOG.md}:379-380); in the
 quadratic model, and conditional on its interfaces, mixing in individually
 squared scores of the other orbit cannot improve the limit, about $5.44209$, of
 the square of the average
 (\path{research/further_2026_10_25/global/individual_square_scope.md}). In the floating-point searches (\emph{screens}) reported for row~14, on the
 first tilted piece the parameter $\ulaw=4/5$ has a slightly smaller chamber
 coefficient (about $5.469937$ against $5.469962$) but a larger layer coefficient
 (about $5.477768$ against $5.445397$), and no global ordering of the two is
 proved (\path{research/further_2026_10_25/global/FIRST_BAND_PARAMETER_SCREEN.md});
 on its last piece the square of the average with $\ulaw=51/40$ has the smaller
 chamber coefficient (about $5.451162$ against $5.456784$ for the individual
 squares with $\ulaw=1$) at a larger layer coefficient (about $4.473194$ against
 $4.430848$), and only the complete closing, which prices both, decides such a
 trade (master \S44;
 \path{research/further_2026_10_25/envelopes/finite_l20_parameter_grid.json},
 \path{research/further_2026_10_25/local/l21_pair_square_tilt.json}). The
 individual squares with $\ulaw=3/4$ already lower the chamber coefficient of the
 uniform root at $2^{35}$, inside the piece from $2^{34}$ to $2^{36}$; splitting
 that piece at $2^{35}$ would need a fresh prime partition, a price replay, a
 native comparison and a complete scalar
 (\path{research/further_2026_10_25/SYNTHESIS.md}:178-184). These are historical
 screens, not the selected laws: rows 15--16 use individual squares from
 $2^{30}$ on six existing pieces, then the coherent $39$-score square with
 $\ulaw=67/40$ and floor $1/100$ from $2^{45}$
 (\cref{def:positive-root,cert:positive-root-caps}).
\item \emph{An interpolation of the matching law} with a bare one-query law, which
 keeps the query and the marginals and moves part of the leading $E$ price from
 skew to symmetric: in a numerical screen with envelopes replaced by their limiting coefficients at the tested parameters the reserved cost rises at the rate of about $0.309$ per unit of the
 moved fraction before the restoration credit, which would subtract about $0.119$ from $\mathcal E$, so it was
 not pursued; this is not a lower bound over all envelopes
 (\path{research/further_2026_10_25/envelopes/E_TRADE_SCREEN.md}).
\item \emph{The late-pool threshold} $2^{31}$ is the best of the $30$ threshold
 pools tested at the parameters of row~12, not an optimum over arbitrary pools or
 other parameters
 (\path{research/continuation_2026_10_24/envelopes/TAIL_THRESHOLD.md}).
\end{itemize}

\begin{remark}[The three depth-index changes]\label{rem:record:item13}
At the October~18 checkpoint three changes of the index were computed: the exact
class-three contraction with least square step $144$, a fixed interval count
$\Kint$, and residues refined by $p$-adic and $2$-adic depths (modulo $p^{128}$
with increments $t^2$, $p^{64}\nmid t$, at the class-three primes below $Q_3$,
modulo $2^{129}$ at $2$, with random conditioning depths). At that checkpoint the
fixed $\Kint$ alone lowered $\Theta$ by about $7.3\cdot10^{-4}$ and the depths by only about $2.8\cdot10^{-5}$. All three changes are now part of the selected construction. The required
extension to nonvanishing square classes is proved in
\cref{lem:kernel-depth-classes,lem:imp:padic-envelope}; the finite-filtration
form of \cref{lem:record:common-mean} supplies the depth-energy identity.
The selected finite moment is certified in \cref{cert:index-moment}.
The historical small numerical gain explains why this change alone does not
provide a substantial exponent improvement.
\end{remark}

\paragraph{Instructive choices in the spectral method.}
\begin{itemize}
\item \emph{Centring at the mean under the actual half marginal $m_p$} leaves
 $\E_\mu[F_\perp\mid x_k]=\kappa_F$ at every vertex, and this constant $\kappa_F$ enters every star and $K_{2,2}$ term; centring at the mean under the half law
 $\mu$ removes these terms exactly, at the price of a second-order cross term
 (\cref{subsec:spec:centring}, \cref{prop:spec2:mucentre}).
\item \emph{Softening:} interpolating each cross constraint
 $\mathbf 1_{\Qsq}$ by $(1-t)/2+t\,\mathbf 1_{\Qsq}$, the derivative in one edge
 isolates $(\chi-\one_{\{0\}})/2$. Its character part has norm $1/(2\sqrt p)$ and its diagonal part norm $1/(2p)$ under uniform normalization; integration
 along a path gives relative cut estimates of order $p^{-1/2}$. They served the
 high range of the $216$-slot design; the final caps come from the two routes of
 \cref{sec:spectral} (\cref{thm:spectral-criterion,thm:record:direct-star}) and,
 at the chamber cut of the middle root, from complete actual stars
 (\cref{subsec:middle-l20,sec:positive-roots}).
\item \emph{Tensor stability} needs no Perron--Frobenius argument or spectral
 gap: top singular vectors of a tensor product can be chosen as products of those of the
 factors, also for non-symmetric kernels. Block bounds in the
 final rows of the table are obtained from one-prime bounds by
 tensoring, using the multiplicativity of chi-square defects
 (\cref{lem:spec:tensorcomp}(f)) and the exactness of the block
 marginals.
\item \emph{The three-block test.} With the first-order four-cycle bound the
 inner $K_{2,2}$ term, bounded through the pure-pair feature-family bound
 $\sum_J\|\Pi_{H_J}F\|^2\le\mathrm{Fr}\,\|F\|^2$, is of order one (about $9$ at
 $2^{44}$), and a three-block test charged it to the pure-pair component; with
 the bound over
 $\mathrm{GL}_2$ it is $O(r)$ (about $4\cdot10^{-6}$), and a two-block test,
 charging it inside the scalar $\bar c_\perp$, suffices.
\item \emph{Class-one refinements.} At $p\equiv1\pmod4$ the contributions of the even edge subsets of the coupling term vanish under a nonsquare dilation, and the internal-edge Paley
 sum $\Sigma_2=\sum_{t,u}\chi(t(t-1)u(u-1)(u-t))$ satisfies $|\Sigma_2|\le2p$, by
 Kloosterman, Gauss and Jacobi sums (numerically $\Sigma_2=2(a^2-b^2)$ for
 $p=a^2+b^2$, $a$ odd). The stored tail caps use neither, only the cruder
 class-one bounds of \cref{lem:spec2:paley-one} and, at the root cuts, the
 parity transfer of \cref{lem:spec2:parity-transfer}.
\item \emph{Restricted criteria} for functions of a vertex subset $A$ supplied
 row envelopes to the $216$-slot completions; their only new feature,
 $\E_\mu[F_\perp\mid x_c]$ non-constant for $c\notin A$, costs outer stars and
 outer $K_{2,2}$ terms. The final completions use full-cut caps.
\end{itemize}

\subsection{Errors to avoid when checking or extending the proof}\label{subsec:obs:pitfalls}

These rules were learned from errors found during the work and from the
reviews; the final chain observes each of them.
\begin{itemize}
\item Never normalize a dominated measure (\cref{rem:cc:mixed}). Work in the
 actual half marginal; an induced half law needs its completion-density
 comparison. Class-three inputs (the common middle law, the Paley counts, the
 class-three caps) are never applied at class-one primes.
\item An upper bound is never a lower bound: the lower endpoint of a certified
 interval enclosing an upper bound for a norm is not a lower bound for the norm,
 which needs its own witness, such as a Rayleigh quotient. Polynomials are
 evaluated at upper prices only because their coefficients are nonnegative
 (\cref{rem:cp:monotone}).
\item Never subtract unrelated upper bounds: a discount enters only as the
 nonnegative coefficient of its class, fixed before the series is rewritten and
 evaluated at the upper prices, and an exponential of an upper price, such as
 the background factor $e^{t\Uinv}$, is subtracted only inside a majorant with
 nonnegative coefficients in the same tag variables
 (\cref{rem:completion-subtraction}); and a jointly certified cap is never
 replaced by a sum of looser bounds.
\item Savings of one degree do not simply add: an exact product can improve the
 exponential envelope but does not automatically add to an existing discount of
 that degree, and of two bounds for one degree the smaller is kept. The indices
 are actual tensor factors, not cells or bands (splitting every factor into
 equal pieces preserves $S$ and $A$ and approaches the envelope). For numbers $b,r$ and $q\ge0$, from $b+q-r$ keep $(b+q-r)_+$, not $(b-r)_++q$; a credit shared by all the orientation signs jointly is established only after summing over the signs; a diagonal added to an upper envelope gives
 a mean credit only if it is in the comparison identity.
\item A variance credit cannot be kept with a smaller envelope: the credit inside
 an envelope class is $c\,g$ with $c\le\kappa$ (\cref{lem:class-credit}).
\item Let $U\ge0$ be a moment (an $L^r$ norm of a fixed
 function under a fixed positive measure) that was priced at
 exponent $a_0$, and let $W\ge0$ be a lower moment of the same
 function under the same measure, not yet priced (\emph{unpriced}).
 Replacing $U$ by $W$ after repricing improves the bound only if
 $W>U/C$, where $C=e^{t_0L}$ with $t_0$ the repricing increment
 $2(a-a_0)$ and $L=570$ the block budget; and $U,W$ must be
 moments for one common positive measure. Local losses
 recur at every induction step and are charged at every step, as $R_I$ is.
\item A larger kernel cutoff or error exponent alone is not a closing theorem.
 The square step, the rounded child scale $\ell$ and its factor $\lambda$ (\cref{lem:record:fibre-scale}(b)) must be kept; a
 same-input scalar inequality does not tensorize across correlated primes, and a
 label shared by several slots does not make them independently averaged.
\item Parameters may be chosen by floating-point search; every inequality used is certified in outward arithmetic
 (\cref{app:certificates}). The integer $\Kint$ is evaluated using the certified
 majorant. Optimality of the square-shift contraction is claimed only within
 its stated class of transitions, and failure of $\Theta<1$ at nearby parameters
 $(a,\Kint)$ concerns that majorant only. Feasibility is not
 monotone in $a$ when a design's conditions are tuned to its exponent.
\item The $\sqrt2$ slip. For $f+Ag=u$, $g+A^*f=v$ with $\|A\|\le\rho<1$ and
 $\|u\|,\|v\|\le U$, the direct-sum norm gives only
 $\max(\|f\|,\|g\|)\le\sqrt2\,U/(1-\rho)$; the bound $U/(1-\rho)$ of the one-way
 repairs needs the equation-by-equation estimates of the proof of
 \cref{lem:tr:repair}, which an earlier proof of the one-way repair skipped.
\item A coupled $0$--$1$ mask is not an operator contraction:
 $\bigl(\begin{smallmatrix}1&1\\1&-1\end{smallmatrix}\bigr)$ has norm $\sqrt2$, and
 zeroing its lower right entry gives $(1+\sqrt5)/2$. So a bound for a complete
 kernel does not justify deleting the collision terms of a grouped extraction,
 and the extraction keeps every special-value and merge allowance (master \S43,
 \texttt{adv:upper:injective-obstruction},
 \path{paper/sections/advance_2026_10_24/upper_interim.tex}:181-207).
\item In pinned pair Grams exclude the pin labels: on $\one_a-\one_b$ at the pins
 $a,b$ the two-pin kernel has the Rayleigh quotient $-q+4q^2-3q^3$, of order $q$,
 while on the mean-zero functions supported off the pins its norm is at most
 $2qr+3q^2$
 (\path{paper/sections/advance_2026_10_24/envelopes_interim.tex}:114-152).
\item A positive diagonal response is not a negative spectral bound, cross
 moments of chamber scores need not be nonnegative, and a bound on a tilted law's
 own marginal is not the fixed-chart bound that the completion needs
 (\path{paper/sections/advance_2026_10_24/local_interim.tex}:207-221, 417-418;
 \cref{sec:positive-roots}).
\item A local improvement is never inserted into an old closing with stale
 dependent quantities: every reference weight, completion inequality and
 assigned price is regenerated, and the gains of different changes interact
 (\path{research/advance_2026_10_24/SYNTHESIS.md}:21-24).
\item Inside the remainder polynomial of a class only the separately proved
 leading and remainder prices may stand: neither a raw-arm price nor a full price
 minus a remainder price may replace a leading price there, since neither bounds
 the leading part of a record change (\cref{rem:completion-subtraction}).
\item An invariant coefficient is not spent twice: a class priced at $\kappa_b$
 on the invariant background keeps only the variance allocation proved for its
 envelope (\cref{lem:class-credit}); and orbit coupons come from an explicit
 decomposition into invariant positive couplings, not from symmetrizing signed
 coefficients (\cref{lem:orbit-coupon}).
\end{itemize}

\subsection{Scope of the reviews}\label{subsec:obs:reviews}

The constructions, proofs and computations were developed in several
parallel lines of work carried out with AI models, an OpenAI Codex model
(gpt-6-astra) and Anthropic's Claude models, partly building on four
earlier research reports handed over between those lines of work. The
initial synthesized text was written with Claude models; the
29~October revision integrates new arguments developed with
ChatGPT Pro Astra~6 and Codex, replacing the corresponding local
estimates and
shortening the historical account. Its exact checks and their
limits are stated in \cref{cert:astra-local}. The reviews are internal reviews by AI models, not external
referee reports. The review (called the record gate in its record) of the October~18 checkpoint
(\path{research/record_gate_2026_10_18/}), with referees of the other model
family, refereed every component of that chain in a mathematics or computation
lens, re-implemented its computations, and passed every component with minor
findings. It checked the research notes and computations, not the text; relied
on the named external theorems; spot-checked without re-running three large
inherited computations (the histogram tree of the upper nonstar classes, the
prime histogram of the middle blocks, the middle spectral grid); and did not
re-referee the earlier sections written with Claude models, whose only review by
the other family was a scoped audit whose three findings the final chain
respects. The components of 24~October that row~11 added (the pure paired law
with its envelopes and restorations, the exact index law with fixed $\Kint$, the
ordered collision bound) were checked by independent re-implementations and
audits (\path{research/push_2026_10_24/}); no referee gate of the other family is
recorded for them. The components that rows 12--14 added (the exact actual laws
and collected channels, the signed whole-block law, the bracket with one fold
copy and the two-cut error, the matrix envelopes and orbit coupons, the complete
actual stars and the direct chamber cap, the dense upper comparisons and joint
rows, the positive roots and their restoration, the assigned invariant price)
were checked by reviewers within the agent team and by independent
re-implementations, recorded with them
(\path{research/continuation_2026_10_24/global_review/CROSS_REVIEW.md},
\path{research/continuation_2026_10_24/upper_arithmetic/LOCAL_PROOF_REVIEW.md},
\path{research/continuation_2026_10_24/upper_arithmetic/GLOBAL_PROOF_REVIEW.md},
\path{research/advance_2026_10_24/envelopes/BETA_CROSS_REVIEW.md},
\path{research/advance_2026_10_24/local/GRAM_BACKGROUND_REVIEW.md},
\path{research/advance_2026_10_24/local/GROUPED_FIRST_STAGE_REVIEW.md},
\path{research/advance_2026_10_24/local/GENERAL_STAR_FRAME_REVIEW.md},
\path{research/further_2026_10_25/REVIEW.md},
\path{research/further_2026_10_25/local/LOCAL_PROOF_REVIEW.md}); their closing
number has an author assembly and an independently written one, which agree on
the $45$ compared components
(\path{research/further_2026_10_25/verification/selected.json}), and these records remain frozen. The combined manuscript of 26 October adds independent reviews by the other model family of
\cref{lem:class-credit,lem:merge-repair,lem:upper-stars}, recorded in
\path{research/merge_2026_10_26/CROSS_FAMILY_REVIEW.md} and
\path{research/merge_2026_10_26/upper_stars_review.md}, together with focused
reviews of the selected root laws and analytic conductor interface. These
are internal reviews, not external referee reports. The merge also records
a focused review of the tagged, signed and auxiliary-price forms of
(CP2)--(CP5), \cref{lem:sharp-first,cor:two-cut-series,lem:record:record-expansion},
and their use in the closing
(\path{research/merge_2026_10_26/comparison_review.md}). This review corrected
two abstract edge cases: degree-class replacement now requires even degree
at least two, and redistribution of upper prices treats zero actual totals
separately. All selected classes have degrees two, four or six and positive
totals, so neither correction changes the bound.

The refined reference-fold cap, the pair-feature matrix at general strength
and the root-normalization prices
(\cref{lem:old-f-law,lem:root-pair-matrix,lem:root-normalization}) have a separate
focused review and independent finite algebra checks in
\path{research/merge_2026_10_26/inherited_local_review.md}. The pair-score
symmetry lemma is explicitly scoped to its two families; coherent-root
admissibility at the selected parameter is
\cref{lem:imp:coherent39-law}. These reviews establish the stated interfaces
under their cited upstream hypotheses, rather than claiming a new
independent proof of every local input or an external review of the entire
manuscript. The weighted star-marginal bound of the archived text is used again here (\cref{lem:wt:stars}), with its density step referred to
\cref{lem:graphs:vertex}. The certificate verifies finite assertions and the
authenticity of their inputs, not the analytic reductions or the external
theorems (\cref{subsec:cert:model}).

The 30~October revision follows two full source-reading passes by a single
Codex reviewer. The principal repairs retain the raw total-variation factor
in the density-sensitive comparison and sum it only after splitting by the
full selected modulus; distinguish fixed-selection asymptotics from the
uniform estimates needed for the complete expansion; and apply the physical
pair normalization once, to raw endpoint coefficients. The revision also
corrects the outer-logarithm remainder in \cref{app:openai-constant}, operator
domains, and several baseline-versus-final numerical attributions. The
selected exponent and frozen numerical construction are unchanged. The
line-coverage record, reproducible checks, findings and remaining limits are
recorded in \path{analysis/proof_repair_review_2026_10_30.json} and its
companion text report. This is an internal adversarial review, not an
independent external referee report or a formal verification of the proof.

\paragraph{Consequences.}
Nothing here is a premise of the main theorem. The numbers quoted are those of
\cref{app:certificates}, \cref{tab:closing-values}, the checkpoints recorded in
\cref{subsec:history:checkpoints}, the research notes named in the text and the
fragments and research notes of master \S\S42--44 named in the text; the
enumerations of \cref{subsec:obs:class-one,subsec:obs:polynomials} are
corroboration only.
